% GENERATED FILE. Edit canonical lessons, then run npm run generate:books.
% canonical-source-hash: fnv1a64:167ee7b2aa123c8b
% canonical-lessons: SA-C449-avegah, SA-C449-vismayate, SA-C449-ascaryam, SA-C449-tosayati, SA-C449-romancah

\chapter{Amazed and Pleased}
\label{ch:sa-amazed-and-pleased}
\hlchaptermodality{449}

\begin{chapteropening}
\textbf{By the end of this chapter:} \emph{I can say what amazed or pleased someone.}

Complete the last lesson of chapter 449: I can say what amazed or pleased someone.
\end{chapteropening}

\section[\texorpdfstring{\sk{आवेगः} (āvegaḥ)}{आवेगः (āvegaḥ)}]{\sk{आवेगः} (āvegaḥ) — agitation; a flurry of feeling}

\label{lesson:SA-C449-avegah}



\begin{quote}
Before the new one: say the Sanskrit for merriment, then the Sanskrit for distress.
\end{quote}

\subsection*{You'll want to know: \sk{आवेगः}}
\textbf{\sk{आवेगः}} — \emph{āvegaḥ} — \textquotedblleft{}agitation; a flurry of feeling\textquotedblright{}.

\begin{quote}
\sk{वार्तया} \sk{मातुः} \sk{आवेगः} \sk{भवति।}
\end{quote}

\begin{quote}
\emph{vārtayā mātuḥ āvegaḥ bhavati.}
\end{quote}

\begin{quote}
The news throws mother into a flurry.
\end{quote}

A masculine noun.

\begin{cousinweb}
\textbf{\sk{आवेगः}} (\emph{āvegaḥ}) is built from \emph{ā} and the root √vij, \textquotedblleft{}to tremble, to start\textquotedblright{}: a quick stirring of the mind; MW gives \textquotedblleft{}hurry or haste produced by excitement\textquotedblright{} and \textquotedblleft{}flurry, agitation\textquotedblright{} (MW).
\end{cousinweb}

\begin{grammarlens}[title={What you've built}]
One more word for describing what you lived and felt.
\end{grammarlens}

\subsection*{Your turn}
\begin{itemize}
\raggedright
  \item \emph{Say it:} \emph{āvegaḥ}
  \item \emph{Say it:} \emph{āvegaḥ}, once more
  \item \emph{Say it:} \emph{vārtayā mātuḥ āvegaḥ bhavati.} — The news throws mother into a flurry.
  \item \emph{Recall:} say the Sanskrit for merriment, then the Sanskrit for distress, then say \emph{āvegaḥ} again
\end{itemize}

\begin{tcolorbox}[breakable,colback=teal!4,colframe=teal!35!black,title={Before you move on}]
What does \sk{आवेगः} mean? (\textquotedblleft{}Agitation; a flurry of feeling\textquotedblright{}.) Say it once more.
\end{tcolorbox}
\section[\texorpdfstring{\sk{विस्मयते} (vismayate)}{विस्मयते (vismayate)}]{\sk{विस्मयते} (vismayate) — is amazed}

\label{lesson:SA-C449-vismayate}



\begin{quote}
Before the new one: say the Sanskrit for distress, then the Sanskrit for a flurry of feeling.
\end{quote}

\subsection*{You'll want to know: \sk{विस्मयते}}
\textbf{\sk{विस्मयते}} — \emph{vismayate} — \textquotedblleft{}is amazed\textquotedblright{}.

\begin{quote}
\sk{रात्रौ} \sk{मीरा} \sk{ताराः} \sk{पश्यति}, \sk{विस्मयते} \sk{च।}
\end{quote}

\begin{quote}
\emph{rātrau mīrā tārāḥ paśyati, vismayate ca.}
\end{quote}

\begin{quote}
At night Mira sees the stars and is amazed.
\end{quote}

Past (she …): \sk{सा} \sk{विस्मितवती}, \emph{sā vismitavatī}.

\begin{cousinweb}
\textbf{\sk{विस्मयते}} (\emph{vismayate}) is built from \emph{vi}, \textquotedblleft{}apart\textquotedblright{}, and the root √smi, \textquotedblleft{}to smile\textquotedblright{}, the verb behind \textbf{\sk{विस्मयः}}; MW gives \textquotedblleft{}to wonder, be surprised or astonished at\textquotedblright{} (MW).
\end{cousinweb}

\begin{grammarlens}[title={What you've built}]
One more word for describing what you lived and felt.
\end{grammarlens}

\subsection*{Your turn}
\begin{itemize}
\raggedright
  \item \emph{Say it:} \emph{vismayate}
  \item \emph{Say it:} \emph{vismayate}, once more
  \item \emph{Say it:} \emph{rātrau mīrā tārāḥ paśyati, vismayate ca.} — At night Mira sees the stars and is amazed.
  \item \emph{Recall:} say the Sanskrit for distress, then the Sanskrit for a flurry of feeling, then say \emph{vismayate} again
\end{itemize}

\begin{tcolorbox}[breakable,colback=teal!4,colframe=teal!35!black,title={Before you move on}]
What does \sk{विस्मयते} mean? (\textquotedblleft{}Is amazed\textquotedblright{}.) Say it once more.
\end{tcolorbox}
\section[\texorpdfstring{\sk{आश्चर्यम्} (āścaryam)}{आश्चर्यम् (āścaryam)}]{\sk{आश्चर्यम्} (āścaryam) — a marvel}

\label{lesson:SA-C449-ascaryam}



\begin{quote}
Before the new one: say the Sanskrit for a flurry of feeling, then the Sanskrit for is amazed.
\end{quote}

\subsection*{You'll want to know: \sk{आश्चर्यम्}}
\textbf{\sk{आश्चर्यम्}} — \emph{āścaryam} — \textquotedblleft{}a marvel\textquotedblright{}.

\begin{quote}
\sk{अहो} \sk{आश्चर्यम्}!
\end{quote}

\begin{quote}
\emph{aho āścaryam!}
\end{quote}

\begin{quote}
Oh, what a marvel!
\end{quote}

A neuter noun.

\begin{cousinweb}
\textbf{\sk{आश्चर्यम्}} (\emph{āścaryam}) MW reports that it is said to come from \emph{ā} and the root √car, \textquotedblleft{}to move\textquotedblright{}, with an \emph{s} inserted (Pāṇini 6.1.147); its first sense is \textquotedblleft{}appearing rarely\textquotedblright{}, and the history is not settled (MW).
\end{cousinweb}

\begin{grammarlens}[title={What you've built}]
One more word for describing what you lived and felt.
\end{grammarlens}

\subsection*{Your turn}
\begin{itemize}
\raggedright
  \item \emph{Say it:} \emph{āścaryam}
  \item \emph{Say it:} \emph{āścaryam}, once more
  \item \emph{Say it:} \emph{aho āścaryam!} — Oh, what a marvel!
  \item \emph{Recall:} say the Sanskrit for a flurry of feeling, then the Sanskrit for is amazed, then say \emph{āścaryam} again
\end{itemize}

\begin{tcolorbox}[breakable,colback=teal!4,colframe=teal!35!black,title={Before you move on}]
What does \sk{आश्चर्यम्} mean? (\textquotedblleft{}A marvel\textquotedblright{}.) Say it once more.
\end{tcolorbox}
\section[\texorpdfstring{\sk{तोषयति} (toṣayati)}{तोषयति (toṣayati)}]{\sk{तोषयति} (toṣayati) — pleases; gratifies}

\label{lesson:SA-C449-tosayati}



\begin{quote}
Before the new one: say the Sanskrit for is amazed, then the Sanskrit for a marvel.
\end{quote}

\subsection*{You'll want to know: \sk{तोषयति}}
\textbf{\sk{तोषयति}} — \emph{toṣayati} — \textquotedblleft{}pleases; gratifies\textquotedblright{}.

\begin{quote}
\sk{पौत्रः} \sk{मधुरेण} \sk{गीतेन} \sk{पितामहं} \sk{तोषयति।}
\end{quote}

\begin{quote}
\emph{pautraḥ madhureṇa gītena pitāmahaṃ toṣayati.}
\end{quote}

\begin{quote}
The grandson pleases his grandfather with a sweet song.
\end{quote}

Past (he …): \sk{सः} \sk{पितामहं} \sk{तोषितवान्}, \emph{saḥ pitāmahaṃ toṣitavān}.

\begin{cousinweb}
\textbf{\sk{तोषयति}} (\emph{toṣayati}) is built on the causative of the root √tuṣ, \textquotedblleft{}to be satisfied, to be pleased\textquotedblright{}, the root of \textbf{\sk{तुष्यति}} and \textbf{\sk{सन्तोषः}}: to make someone pleased; MW gives the causative the senses \textquotedblleft{}to satisfy, please, appease\textquotedblright{} (MW).
\end{cousinweb}

\begin{grammarlens}[title={What you've built}]
One more word for describing what you lived and felt.
\end{grammarlens}

\subsection*{Your turn}
\begin{itemize}
\raggedright
  \item \emph{Say it:} \emph{toṣayati}
  \item \emph{Say it:} \emph{toṣayati}, once more
  \item \emph{Say it:} \emph{pautraḥ madhureṇa gītena pitāmahaṃ toṣayati.} — The grandson pleases his grandfather with a sweet song.
  \item \emph{Recall:} say the Sanskrit for is amazed, then the Sanskrit for a marvel, then say \emph{toṣayati} again
\end{itemize}

\begin{tcolorbox}[breakable,colback=teal!4,colframe=teal!35!black,title={Before you move on}]
What does \sk{तोषयति} mean? (\textquotedblleft{}Pleases; gratifies\textquotedblright{}.) Say it once more.
\end{tcolorbox}
\section[\texorpdfstring{\sk{रोमाञ्चः} (romāñcaḥ)}{रोमाञ्चः (romāñcaḥ)}]{\sk{रोमाञ्चः} (romāñcaḥ) — a thrill; the hair standing on end}

\label{lesson:SA-C449-romancah}



\begin{quote}
Before the new one: say the Sanskrit for a marvel, then the Sanskrit for pleases.
\end{quote}

\subsection*{You'll want to know: \sk{रोमाञ्चः}}
\textbf{\sk{रोमाञ्चः}} — \emph{romāñcaḥ} — \textquotedblleft{}a thrill; the hair standing on end\textquotedblright{}.

\begin{quote}
\sk{सिंहस्य} \sk{गर्जनेन} \sk{मम} \sk{रोमाञ्चः} \sk{भवति।}
\end{quote}

\begin{quote}
\emph{siṃhasya garjanena mama romāñcaḥ bhavati.}
\end{quote}

\begin{quote}
The lion's roar gives me a thrill.
\end{quote}

A masculine noun.

\begin{cousinweb}
\textbf{\sk{रोमाञ्चः}} (\emph{romāñcaḥ}) is built from \emph{roman}, \textquotedblleft{}the hair of the body\textquotedblright{}, and \emph{añca}, from the root √añc, \textquotedblleft{}to bend, to rise\textquotedblright{}; MW gives \textquotedblleft{}thrill of the hair\textquotedblright{} (MW).
\end{cousinweb}

\begin{grammarlens}[title={What you've built}]
One more word for describing what you lived and felt.
\end{grammarlens}

\subsection*{Your turn}
\begin{itemize}
\raggedright
  \item \emph{Say it:} \emph{romāñcaḥ}
  \item \emph{Say it:} \emph{romāñcaḥ}, once more
  \item \emph{Say it:} \emph{siṃhasya garjanena mama romāñcaḥ bhavati.} — The lion's roar gives me a thrill.
  \item \emph{Recall:} say the Sanskrit for a marvel, then the Sanskrit for pleases, then say \emph{romāñcaḥ} again
\end{itemize}

\begin{tcolorbox}[breakable,colback=teal!4,colframe=teal!35!black,title={Before you move on}]
What does \sk{रोमाञ्चः} mean? (\textquotedblleft{}A thrill; the hair standing on end\textquotedblright{}.) Say it once more.
\end{tcolorbox}
